\section{Handing a Field Over}
\label{sec:ch17-handing-a-field-over}

\index{override@\code{"@override}!the migration it exists for|(}%
Section~\ref{sec:ch17-the-field-nobody-approved} used \code{@override} to do
something nobody should do, which is not an argument against the directive.
It is the mechanism the specification offers for the thing every federated
graph eventually needs, moving a field from the service that has it to the
service that should. Apollo's reference says exactly
that~\autocite{apollodirectives}: it \enquote{indicates that an object field is
now resolved by this subgraph instead of another subgraph where it is also
defined}. Used in order, it works.

\index{override@\code{"@override}!the order the steps have to go in}%
The order is the whole of it, and none of it is enforced. The destination
service implements the field and deploys, so that something can answer it.
Then its document publishes the declaration with \code{@override(from:)} on it,
which is the step composition approves. Then the origin service removes its
copy, at whatever later date suits it, because until then the field has two
declarations and only one route. Composition sits in the middle of that
sequence and can see one step of it. Whether the first ever happened is the
question section~\ref{sec:ch17-the-field-nobody-approved} answered with a
request, and the answer took a deployment to arrive.

\subsection{The Restriction That Is Enforced}

\index{override@\code{"@override}!only one subgraph may take a field}%
One rule around the directive is checked, and it is the rule that would
otherwise leave the composer with the problem
section~\ref{sec:ch17-the-one-thing-it-asks-permission-for} is about. Apollo
states it as \enquote{only one subgraph can \code{@override} any given field.
If multiple subgraphs attempt to \code{@override} the same field, a
composition error occurs}~\autocite{apollodirectives}. Put the same
\code{@override(from: "sessions")} on \code{Session.abstract} in both the
Ratings and the Search documents and \code{wgc} refuses, saying that the
directive must be declared on one single instance of a field and naming both
subgraphs.

\index{override@\code{"@override}!which is a routing rule again}%
Which is the same rule as the shareable one wearing different words. Two
services taking a field leaves two candidates and no way to choose, so it
stops. One service taking a field leaves one, so it does not.

\subsection{The Half This Composer Does Not Have}

\index{override@\code{"@override}!progressive, and who ships it}%
Apollo's reference carries a second argument on the directive, \code{label},
optional, whose purpose is to split traffic between the two services while a
migration is in flight~\autocite{apollodirectives}. Hot Chocolate ships it:
\code{OverrideAttribute} in \code{HotChocolate.ApolloFederation} carries a
\code{Label} beside its \code{From}, and the doc comment on the attribute
points at Apollo's own page for what a percentage label means. So a .NET
subgraph author can write it, and the document that service exports will carry
it.

\index{override@\code{"@override}!the composer refuses the label}%
Put the same directive and the same definition into the Ratings document by
hand, which is what that service would have exported, and compose the four:

\begin{minted}[fontsize=\footnotesize,breakanywhere]{text}
The subgraph "ratings" could not be federated for the following reason:
The 1st instance of the directive "@override" declared on coordinates "Session.abstract" is invalid for the following
reason:
 The definition for "@override" does not define the following argument that is provided: "label".
\end{minted}

\index{override@\code{"@override}!the document defines it and is overruled}%
The document does define it, and this is the definition it carries:

\begin{graphqlsdl}
directive @override(from: String!, label: String) on FIELD_DEFINITION
\end{graphqlsdl}

That is what the composer declines to read. It holds a definition of its own,
that definition has one argument, and it substitutes it. This is
the fourth time in this book that the two libraries have disagreed about the
shape of a directive, after the cost weight in chapter~\ref{ch:composition},
the \code{@provides} precondition in
chapter~\ref{ch:second-third-subgraph} and the semantic-non-null levels
argument in chapter~\ref{ch:null-blast-radius}.

\index{override@\code{"@override}!Cosmo's answer is a different mechanism}%
This one has a stated reason rather than a drift. Jens Neuse, a co-founder of
WunderGraph, describes progressive override as Apollo's feature, one that
\enquote{allows you to annotate a field with the \code{@override} directive
and an additional \code{label} argument to gradually move ownership of a field
from one Subgraph to another}, and puts Cosmo's own answer somewhere
else~\autocite{neuse2024featureflags}: a feature flag over a whole alternative
subgraph, enabled for a fraction of traffic, rather than an argument on a
field. He is describing the product he sells, and he is describing it
accurately in both directions. What it means for a reader on this stack is
plain enough: the incremental migration Apollo documents is not available, and
the field changes hands in one deployment.

\subsection{Two Ways to Leave It Wrong}

\index{override@\code{"@override}!a mistyped subgraph name}%
The name in \code{from} is the subgraph's, the one \code{graph.yaml} gives it,
and getting it wrong is not reported as getting it wrong.
\code{@override(from: "session")}, singular, is an easy thing to type, and
what it does is nothing at all: the directive fails to apply and what comes
back is the refusal from
section~\ref{sec:ch17-the-one-thing-it-asks-permission-for}, that two subgraphs
define \code{abstract} and neither declares it shareable. The message names
neither the directive nor the string that was typed. A reader who has just
written an override and gets a message about \code{@shareable} has to work out
for themselves that those are the same event.

\index{override@\code{"@override}!outlives the field it took}%
And at the far end of a migration the directive stops meaning anything and
nothing says so. Delete \code{abstract} from the Sessions document, which is
the last step and the correct one, and the Ratings copy still reads
\code{@override(from: "sessions")}, pointing at a declaration that no longer
exists in a service that no longer has the field. That composes at exit zero.
The directive is now a comment, it looks like configuration, and the next
person to read that document will believe there is a migration in progress.

\subsection{The Directive This Book Does Not Ship}

\index{provides@\code{"@provides}!why the book has none}%
Chapter~\ref{ch:second-third-subgraph} met \code{@provides} and shipped none,
and the reason belongs here. The directive is a routing optimization: it tells
the router that this particular path through the graph can produce some of
another service's fields locally, so the entity fetch can be skipped. Its
preconditions are the problem. The provided field has to be declared
\code{@external} where it is provided and \code{@shareable} in at least one
other subgraph that defines it, which means the service that owns none of the
column is nonetheless declaring the column, and some other service has agreed
in writing that a second answer to that field is fine.

\index{provides@\code{"@provides}!a copy of somebody else's column}%
That is a copy of somebody else's data living in your service, with a schema
declaration to match, so that a fetch can be avoided. It is the exact
duplication chapter~\ref{ch:second-third-subgraph} removed from the Sessions
service when it moved the speaker rows out. I would rather pay the fetch. A
graph where three services carry declarations of each other's columns is a
graph where the question this chapter opens with, who owns this field, has no
answer left at all.
\index{override@\code{"@override}!the migration it exists for|)}%
